\documentclass[a4paper,11pt]{article}
\usepackage[T2A]{fontenc}
\usepackage[utf8]{inputenc}
\usepackage[russian]{babel}
\usepackage{amsmath,amssymb,mathtools}
\usepackage[margin=21mm,top=23mm,bottom=24mm]{geometry}
\usepackage{array,booktabs,longtable}
\usepackage{xcolor}
\usepackage{microtype}
\usepackage{tikz}
\usepackage{hyperref}
\usepackage{fancyhdr}
\usepackage{enumitem}
\definecolor{accent}{HTML}{087D86}
\definecolor{muted}{HTML}{51565B}
\definecolor{rednote}{HTML}{AB293B}
\hypersetup{hidelinks,pdfauthor={Лёвин Артём Александрович},pdftitle={Ревью домашней работы — Екатерина Скелли}}
\setlength{\parindent}{0pt}
\setlength{\parskip}{6pt}
\setlist[itemize]{leftmargin=*,topsep=2pt,itemsep=2pt}
\pagestyle{fancy}
\fancyhf{}
\fancyhead[L]{\small\textcolor{muted}{Физика · Плавление и единицы измерения}}
\fancyhead[R]{\small\textcolor{muted}{Екатерина Скелли}}
\fancyfoot[C]{\small\thepage}
\renewcommand{\headrulewidth}{0.3pt}
\setlength{\headheight}{15pt}
\newcommand{\TaskTitle}[2]{\phantomsection\label{task:#1}\section*{Задание №#1. #2}}
\newcommand{\Field}[1]{\par\vspace{3pt}{\bfseries #1}\par\vspace{-2pt}}
\newcommand{\ErrorLabel}{\textcolor{rednote}{\textbf{Ошибка:}}}
\newcommand{\IdeaLabel}{\textcolor{accent}{\textbf{Ключевая идея:}}}
\newcommand{\ResultLabel}{\textcolor{accent}{\textbf{Итог:}}}
\begin{document}
\thispagestyle{empty}
\begin{center}
\vspace*{28mm}
{\small\color{accent}\textsc{Физика}}\par
\vspace{17mm}
{\Huge Ревью домашней работы}\par
\vspace{10mm}
{\Large\color{accent} Плавление и количество теплоты}\par
\vspace{8mm}
{\large Перевод единиц и проверка вычислений}\par
\vfill
{\Large Екатерина Скелли}\par
\vspace{12mm}
15 ноября 2026 г.\par
\vspace{17mm}
{\small Лёвин Артём Александрович}\par
{\small эксклюзивно для Екатерины Скелли}\par
\vspace{16mm}
\begin{tikzpicture}
\draw[accent,line width=0.8pt] (-4,0)..controls(-2.3,0.9)and(-1.3,-0.6)..(0,0)
 ..controls(1.4,0.7)and(2.7,-0.9)..(4,0);
\end{tikzpicture}
\end{center}
\clearpage
\section*{Сводная таблица}
\small
\renewcommand{\arraystretch}{1.38}
\begin{longtable}{@{}p{1.3cm}p{2.15cm}p{5cm}p{2.55cm}p{2.55cm}@{}}
\toprule
\textbf{№} & \textbf{Тема} & \textbf{Суть затруднения} & \textbf{Ваш ответ} & \textbf{Верный ответ}\\
\midrule
\endfirsthead
\toprule
\textbf{№} & \textbf{Тема} & \textbf{Суть затруднения} & \textbf{Ваш ответ} & \textbf{Верный ответ}\\
\midrule
\endhead
\hyperref[task:3]{3} & Плавление & При подстановке удельная теплота увеличена в десять раз. & 100 кДж & 10 кДж\\
\bottomrule
\end{longtable}
\normalsize
\medskip
\textcolor{muted}{В таблице оставлено только задание, требующее исправления. Номер соответствует исходному листу.}
\clearpage
\thispagestyle{empty}
\vspace*{0.34\textheight}
\begin{center}{\LARGE\textcolor{accent}{Разбор ошибок}}\end{center}
\clearpage
\TaskTitle{3}{Плавление свинца}
\Field{Текст задания}
Свинец массой 400~г нагрет до температуры плавления. Определите количество теплоты для полного плавления, если удельная теплота плавления равна 25~кДж/кг.
\Field{Ваше решение и ответ}
Вы верно перевели массу в килограммы и записали значение удельной теплоты в СИ:
\[
m=0{,}4\ \text{кг},\qquad \lambda=2{,}5\cdot10^4\ \text{Дж/кг},\qquad Q=\lambda m.
\]
Далее в численной подстановке появляется запись
\[
Q=25\cdot10^4\cdot0{,}4=100000\ \text{Дж}.
\]
Записанный ответ: 100~кДж.
\Field{Где именно возникает ошибка}
\textbf{Последний правильный шаг:} $\lambda=2{,}5\cdot10^4$~Дж/кг. \textbf{Первый неправильный шаг:} вместо $2{,}5\cdot10^4$ подставлено $25\cdot10^4$.
\Field{Почему это неверно}
\ErrorLabel\ При переходе от $2{,}5$ к $25$ коэффициент увеличился в десять раз, но показатель степени остался прежним. Поэтому значение удельной теплоты стало в десять раз больше исходного. Правило перевода таково: $1$~кДж = $10^3$~Дж. Следовательно,
\[
25\ \text{кДж/кг}=25\cdot10^3\ \text{Дж/кг}=2{,}5\cdot10^4\ \text{Дж/кг}.
\]
\Field{Ключевая идея и исправление}
\IdeaLabel\ Сохраните правильный показатель степени при переходе к подстановке. Нужно исправить только коэффициент:
\[
Q=\underbrace{2{,}5\cdot10^4}_{\lambda}\cdot\underbrace{0{,}4}_{m}.
\]
\Field{Полное правильное решение}
При плавлении тела, уже нагретого до температуры плавления, используется $Q=\lambda m$. Нагревание до температуры плавления добавлять не нужно. Переводим данные в СИ:
\[
m=400\ \text{г}=0{,}4\ \text{кг},\qquad
\lambda=25\ \text{кДж/кг}=25000\ \text{Дж/кг}.
\]
Подставляем:
\[
Q=25000\cdot0{,}4=10000\ \text{Дж}=10\ \text{кДж}.
\]
\Field{Проверка}
Поскольку для каждого килограмма нужно 25~кДж, для $0{,}4$~кг потребуется $0{,}4$ от этой величины: $25\cdot0{,}4=10$~кДж. Результат согласован с исходными данными.
\Field{Итог}
\ResultLabel\ $Q=10000$~Дж $=10$~кДж.
\Field{Задача на закрепление}
Олово массой 600~г уже находится при температуре плавления. Найдите количество теплоты для полного плавления при $\lambda=60$~кДж/кг.
\clearpage
\end{document}